\begin{conf-abstract}
{Tracing Hidden Connections: Multi-Tracer Evidence of Managed and Natural Controls on Groundwater–Surface Water Contributions to the Eerste River Estuary, South Africa}
{Welham, A.\textsuperscript{a,b}, van Rooyen, J.\textsuperscript{a,b,c}, Watson, A.\textsuperscript{c}, Roychoudhury, A.\textsuperscript{a}, Chow, R.\textsuperscript{a,d}}
{(a) Stellenbosch University (SU)\\
(b) Swiss Federal Institute of Aquatic Science and Technology (Eawag)\\
(c) School of Climate Studies, SU\\
(d) Wageningen University \& Research.}
{Estuarine wetlands in Mediterranean Climates are highly sensitive to variations in water availability and quality, which influence their ecological services. The Eerste River Estuary in the Western Cape, South Africa, is vulnerable to the combined pressures of recharge variability, urbanisation, and water transfer schemes. In this data-limited, urbanised catchment, a multi-tracer approach was applied to investigate the spatio-temporal contributions of water sources (rain, surface water, and shallow and deep groundwater) to the estuary. Major ions, stable isotopes (\ce{\mathrm{\delta}^{18}O}, \ce{\mathrm{\delta}^{2}H}), tritium, and dissolved gases (\ce{N2}, \ce{O2}, \ce{CH4}, \ce{CO}, \ce{He}, \ce{Ne}, \ce{Ar}, \ce{Kr}, \ce{Xe}) were analysed from 17 rainwater, 83 surface water, and 33 groundwater (shallow and deep aquifers) sites between 2022 and 2024 dry and wet seasons. Hydrochemical facies evolve from the freshwater Na–Cl type in the headwater recharge to \ce{Na-Cl-Ca-HCO3} type near the estuary, reflecting a spatial gradient controlled by precipitation, evaporation and anthropogenic influences, with greatest fluctuations in the estuary chemistry. Isotopic data indicates evaporation losses of the estuary are buffered by a dual-water system: short-term variability is dominated by upstream surface water transfer schemes, while longer-term resilience depends on groundwater of contrasting ages and origins. Groundwater tritium levels are uniformly low ($<$0.5 TU), indicating predominantly pre-modern recharge and long residence times, though some sites (0.5–1.7 TU) suggest minor mixing with modern water. The presence of crustal-derived helium (\ce{^3He}/\ce{^4He} $R/Ra$ = 0 – 0.04) and \ce{^4He}-excess air further supports the dominance of old groundwater with limited modern recharge. These findings show that although the estuary is buffered by managed inflows, it remains vulnerable to climatic variability and abstraction pressure. The multi-tracer framework has refined the conceptual understanding of complex hydrodynamic systems and guides for sustainable monitoring and management of estuarine wetlands.
}
\end{conf-abstract}
